\section{Academic Publications}
	
	\cvitem{}{\begin{center} \textcolor{magenta}{\textcolor{magenta}{Journal Papers}}\end{center}}

    \cvitem{J13}{Cotton, K., de Vries K., Holzapfel A., Berglund K., \textcolor{magenta}{Tatar, K.}. \textit{Revising MIR Research Practices for Singing Data Collection}. Submitted to AI and Society, Springer Nature.}

    \cvitem{J12}{Liu X., Cotton K., \textcolor{magenta}{Tatar, K.} \textit{Revisiting a Decade of Computational Game Creativity}. Submitted to International Journal on Human Computer Interaction, Taylor and Francis.}

    \cvitem{J11}{*Fröhlich M. M., Cotton K., \textcolor{magenta}{Tatar, K.} \textit{Artificial Intelligence: Perspectives on a Caesura in Art, its Creation and Consumption}. Submitted to Leonardo, MIT Press.}

    \cvitem{J10}{Madaghiele V., Lund L., Holzer D., Kelkar T., \textcolor{magenta}{Tatar, K.}, and Holzapfel A. \textit{Expanding the machine: notating generative synthesis with a state-based representation and an interactive timbre space}. Submitted to and passed first round of revisions at Organised Sound, Cambridge Press.}

    \cvitem{J9}{Zappi, V.,\textcolor{magenta}{Tatar, K}. (2025). \textit{Neural audio instruments: Epistemological and phenomenological perspectives on musical embodiment of deep learning}. Frontiers in Computer Science, 7. \url{https://doi.org/10.3389/fcomp.2025.1575168}}

    \cvitem{J8}{Cotton, K., Kaila, A.K.,Jääskeläinen, P., Holzapel, A., \textcolor{magenta}{Tatar, K}. (2025).\textit{Imploding between the Facts and Concerns of Musical-AI: A Multidisciplinary Method for Analysing Human-AI Musical Interaction.} Humanities and Social Sciences Communications, 12(1), 1–20. \url{https://doi.org/10.1057/s41599-025-04533-4}}

    \cvitem{J7}{*\textcolor{magenta}{Tatar, K.}, Ericson P., Cotton K., Núñez del Prado P. T., Batlle-Roca R., Cabrero-Daniel, B, Ljungblad S., Diapoulis G., Hussain J. (2024) \textit{A Shift In Artistic Practices through Artificial Intelligence}. Leonardo Journal, MIT Press.}
 
	\cvitem{J6}{Strauss, L., \textcolor{magenta}{Tatar, K.}, Nuro, S. (2021). Iterative Design Processes and Soma-based Practices in \textit{instance}. Organised Sound, Special Issue on "Collective and Networked Sound Practices.''. Cambridge Press. \url{https://doi.org/10.1017/S1355771821000479}}
		
	\cvitem{J5}{\textcolor{magenta}{Tatar, K.}, Bisig, D., \& Pasquier, P. Latent Timbre Synthesis: Audio-based variational auto-encoders for music composition and sound design applications. The Special Issue of \href{https://www.springer.com/journal/521}{Neural Computing and Applications}: ``Networks in Art, Sound and Design.'' Springer. \url{https://doi.org/10.1007/s00521-020-05424-2}}
	
	\cvitem{J4}{**\textcolor{magenta}{Tatar, K.}. (2020). Initial Remarks on Analyzing Acousmatic Music from the Perspective of Multi-agents. Array, special issue on Agency. International Computer Music Association. http://dx.doi.org/10.25532/OPARA-46}
	
	\cvitem{J3}{*\textcolor{magenta}{Tatar, K.}, Prpa M., \& Pasquier, P. (2019). A Virtual Reality Art Piece with a Musical Agent guided by Respiratory Interaction: \textit{Respire}. \textit{Leonardo Music Journal} Vol. 29, p. 19--24. \textcolor{magenta}{Cover of the issue}. \url{https://doi.org/10.1162/lmj_a_01057} }
	
	\cvitem{J2}{\textcolor{magenta}{Tatar, K.}, \& Pasquier, P. (2018). Musical agents: A typology and state of the art towards Musical Metacreation. \textit{Journal of New Music Research}, 48(1), 56–-105. \url{https://doi.org/10.1080/09298215.2018.1511736}.}
	
	\cvitem{J1}{\textcolor{magenta}{Tatar, K.}, Macret, M., \& Pasquier, P. (2016). Automatic Synthesizer Preset Generation with PresetGen. \textit{Journal of New Music Research}, 45(2), 124--144. \url{http://doi.org/10.1080/09298215.2016.1175481}}
	
	\cvitem{}{\begin{center} \textcolor{magenta}{\textcolor{magenta}{Conference Papers}}\end{center}}

    \cvitem{C20}{Liu, X. and Tatar, K. \textit{AEGIS: Authentic Edge Growth In Sparsity for Link Prediction in Edge-Sparse Bipartite Knowledge Graphs}. Submitted to ICLR 2026.}
 
    \cvitem{C19}{Chen. J., \textcolor{magenta}{Tatar, K.}, Zappi, V. (2024). \textit{A Deep Learning Framework for Musical Acoustics Simulations}. In proceedings of AI Music Creativity 2024.}

    \cvitem{C18}{Cotton, K., \textcolor{magenta}{Tata,r K.}. (2024). \textit{Sounding out extra-normal AI voice: Non-normative musical engagements with normative AI voice and speech technologies}. In proceedings of AI Music Creativity 2024.}

    \cvitem{C17}{Caravati. M, \textcolor{magenta}{Tatar, K.}. (2024). \textit{Interfacing ErgoJr with Creative Coding Platforms}. In proceedings of the 9th International ACM Conference on Movement and Computing.}

    \cvitem{C16}{Cotton. K.,de Vries, K., and \textcolor{magenta}{Tatar, K.}. (2024). \textit{Singing for the Missing: Bringing the Body Back to AI Voice and Speech Technologies}. In proceedings of the 9th International ACM Conference on Movement and Computing.}

    \cvitem{C15}{Cotton. K.,\textcolor{magenta}{Tatar K.} (2024).\textit{ Caring Trouble and Musical AI: Considerations towards a Feminist Musical AI}. In Proceedings of AI Music Creativity Conference 2023. Received the highest ranking (5/5) from all three reviewers. \url{https://aimc2023.pubpub.org/pub/zwjy371l}}

    \cvitem{C14}{\textcolor{magenta}{Tatar K.},  Cotton, K., Bisig, D. (2023). \textit{Sound Design Strategies for Latent Audio Space Explorations Using Deep Learning Architectures}. In Proceedings of Sound and Music Computing 2023.}
 
	\cvitem{C13}{Obaid M., \textcolor{magenta}{Tatar K.}, Wiberg M., Said A., Rost M., Weilenmann A., Johal W., Eyssel F. (2022). \textit{Social Drones for Health and Well-being}. In Adjunct Proceedings of the 2022 Nordic Human-Computer Interaction (Nordic-CHI 2022) Conference.}
 
	\cvitem{C12}{Zappi V,. \textcolor{magenta}{Tatar, K.} (2022) \href{https://embedded-ai-for-nime.github.io/}{\textit{The Neuralacoustics Project: Exploring Deep-Learning for Lightweight Numerical Modeling Synthesis}}. Victor Zappi and K{\i}van\c{c} Tatar. In \textit{Embedded AI for NIME: Challenges and Opportunities} Workshop at New Interfaces for Musical Expression (NIME) 2022.}
	
    \cvitem{C11}{Diapoulis G., Zannos I., \textcolor{magenta}{Tatar, K.}, Dahlstedt P. (2022). Bottom-up live coding: Analysis of continuous interactions towards predicting programming behaviours. In Proceedings of New Interfaces for Music Expression Conference 2022.}
	
	\cvitem{C10}{Bisig D., \textcolor{magenta}{Tatar, K.} (2021). Raw Music from Free Movements: Early Experiments in Using Machine Learning to Create Raw Audio from Dance Movements. Accepted to AI Music Creativity Conference 2021 (AIMC 2021).}
	
	\cvitem{C9}{Boerson R., Liu-Rosenbaum A., \textcolor{magenta}{Tatar K.}, Pasquier P. (2020). Chatterbox: an interactive system of gibberish agents. Accepted to the International Symposium of Electronic Arts (ISEA 2020)}
	
	\cvitem{C8}{\textcolor{magenta}{Tatar K.}, Pasquier P., Siu R. (2019). Audio-based Musical Artificial Intelligence and Audio-Reactive Visual Agents in Revive. In Proceedings of the International Computer Music Conference and New York City Electroacoustic Music Festival 2019 (ICMC-NYCEMF 2019)}
	
	\cvitem{C7}{Fan J., Thorogood M., \textcolor{magenta}{Tatar K.}, Pasquier P. (2018). Quantitative Analysis of the Impact of Mixing on Perceived Emotion of Soundscape Recordings. In Proceedings of Sound and Music Computing (SMC 2018)}
	
	\cvitem{C6}{**\textcolor{magenta}{Tatar K.}, Pasquier P., \& Siu R. (2018). REVIVE: An Audio-Visual Performance with Musical and Visual Artificial Intelligence Agents. 2018 CHI Conference on Human Factors in Computing Systems, Montreal, QC, Canada ACM 978-1-4503-5621-3/18/04. \url{https://doi.org/10.1145/3170427.3177771}}
	
	\cvitem{C5}{**Prpa M., \textcolor{magenta}{Tatar K.}, Schiphorst T., \& Pasquier P. (2018). Respire: A Breath Away from the Experience in Virtual Environment. 2018 CHI Conference on Human Factors in Computing Systems, Montreal, QC, Canada ACM 978-1-4503-5621-3/18/04. \url{https://doi.org/10.1145/3170427.3180282}}
	
	\cvitem{C4}{ Prpa M., \textcolor{magenta}{Tatar K.}, Fran\c{c}oise J., Riecke B., Schiphorts T., Pasquier P. (2018). Attending to Breath: Exploring How the Cues in a Virtual Environment Guide the Attention to Breath and Shape the Quality of Experience to Support Mindfulness. In Proceedings of the 2018 Designing Interactive Systems Conference (pp. 71-84). ACM Press. \url{https://doi.org/10.1145/3196709.3196765}} 
	
	\cvitem{C3}{\textcolor{magenta}{Tatar, K.} \& Pasquier, P. (2017). MASOM: A Musical Agent Architecture based on Self-Organizing Maps, Affective Computing, and Variable Markov Models. In \textit{Proceedings of the 5th International Workshop on Musical Metacreation} (MuMe 2017).}
	
	\cvitem{C2}{Fan, J., \textcolor{magenta}{Tatar, K.}, Thorogood, M., , \& Pasquier, P. (2017). Ranking Based Experimental Music Emotion Recognition. In \textit{Proceedings of the 18th International Society for Music Information Retrieval Conference} (ISMIR 2017).}
	
	\cvitem{C1}{Prpa, M., \textcolor{magenta}{Tatar, K.}, Riecke, B. E., \& Pasquier, P. (2017). The Pulse Breath Water System: Exploring Breathing as an Embodied Interaction for Enhancing the Affective Potential of Virtual Reality. In \textit{Virtual, Augmented and Mixed Reality, 9th International Conference, VAMR 2017, Held as Part of HCI International 2017, Proceedings}. Vancouver: Springer}
	
	\cvitem{}{\begin{center} \textcolor{magenta}{\textcolor{magenta}{Arxiv}}\end{center}}

    \cvitem{AX2}{Dignum V., Casey D., Cerratto-Pargman T., Dignum F., Fantasia V., Formark B., Hammarfelt B., Holmberg G., Holzapfel A., Larsson S., Lagerkvist A., Lakemond N., Lindgren H., Lorig F., Marusic A., Rahm L., Razmetaeva Y., Sikström S., \textcolor{magenta}{Tatar K.}, Tucker J. (Submitted). \textit{On the importance of AI research beyond disciplines.} \textcolor{magenta}{Submitted to \href{https://www.nature.com/ncomms/}{Nature Communications}}. Dignum V. is the first author, and the remaining authors share equal contributions. Preprint: \url{https://arxiv.org/abs/2302.06655}}
 
	\cvitem{AX1}{\textcolor{magenta}{Tatar K.}, Bisig D., \& Pasquier, P. (2020). Introducing Latent Timbre Synthesis. \url{https://arxiv.org/abs/2006.00408}}
	\cvitem{}{\begin{center} \textcolor{magenta}{\textcolor{magenta}{Doctoral Thesis}}\end{center}}
	\cvitem{TH1}{Tatar, K. (2019). Musical agents based on self-organizing maps for audio applications [Thesis, Communication, Art \& Technology: School of Interactive Arts and Technology].\url{http://summit.sfu.ca/item/19665}}
	\cvitem{}{}
	\cvitem{}{\small All papers are double-blind peer-reviewed unless otherwise is indicated.}
	\cvitem{}{\small *These papers are single-blind peer-reviewed. The reviewers had access to the list of authors.}
	\cvitem{}{\small **We are invited to submit a paper, and the article is reviewed by the editors.}